% ---------------------------
% Universidad de Costa Rica
% Facultad de Ingeniería
%
% JORNADAS DE INVESTIGACIÓN
% ---------------------------

% Tipo de documento
\documentclass{jornadas}

% A. Paquetes especiales
% Paquetes y definiciones que no están incluidos 
% en la clase jornadas.cls o que son 
% propios del artículo.
\input{configuracion.tex}

% B. Título del artículo
\titulo{Nam sed sem accumsan, convallis leo eget, commodo ipsum, ut id lectus accumsan est fermentum sollicitudin}

% Título traducido al inglés
\thetitle{Lorem ipsum dolor sit amet, consectetur adipiscing elit, integer sollicitudin magna quis velit fermentum, vitae elementum mauris scelerisque}

% B. Autor(a)(es) del artículo
% Todos los nombres incluyen dos apellidos,
% acentos y signos de puntuación apropiados.
\autor{Nombre completo}
{Labor, unidad académica}
{Institución}
{Correo institucional}
{Número ORCID}

\autor{Juan Rafael Mora Porras}
{Presidente, Poder Ejecutivo}
{Gobierno de Costa Rica}
{juan.mora@costarica.com}
{0123456789}

% Agregar tantos autores como sea necesario con el comando \autor{nombre}{posición}{institución}{correo}{orcid}

% Palabras clave
\palabras{Amortiguamiento, control estructural, espectro de densidad de potencia, modelado dinámico, respuesta en frecuencia}

% C. Bibliografía
% Documento con la bibliografía de este
% artículo, en formato bibtex.
\addbibresource{referencias.bib}

%%%%%%%%%%%%%%%%
\begin{document}
%%%%%%%%%%%%%%%%

%%%%%%%%
\titulos

%%%%%%%%%%%%%%%
\listadeautores

%%%%%%%%%%%%%%
\palabrasclave

% Iniciar numeración de líneas
\linenumbers

%%%%%%%%
\section{Resumen extendido}
%%%%%%%%

El manuscrito a presentar debe ser un resumen extendido de un máximo de 1000 palabras (esta extensión no incluye las figuras, tablas), debe tener una redacción clara y concisa, organizado en una sola columna y dividido en secciones. Las páginas y las líneas deben tener numeración. El documento debe estar escrito con letra Times New Roman de 11 puntos, a doble espacio en tamaño carta\footnote{Este formato debe ser igual para todos los autores o las autoras del artículo. Asimismo, las notas al pie sólo deben utilizarse en caso de ser necesario y seguir este formato. Se debe asegurar también que el número de la nota corresponda con el que señaló en el texto.}. Se sugiere que el contenido de este resumen tenga lo siguiente:  problema o asunto a tratar, pregunta de la investigación o teorías aplicadas, resultados y conclusiones.

El Manual de estilo para citación y referencias que se debe utilizar es el IEEE, puede ser consultado en el siguiente enlance: \url{https://ieeeauthorcenter.ieee.org/wp-content/uploads/IEEE-Reference-Guide.pdf}.

%%%%%%%%
\section{Consideraciones generales}
%%%%%%%%

%%%%%%%%%%%
\subsection{Autorización de uso de materiales}

Es responsabilidad del autor o de la autora obtener el permiso escrito correspondiente en el caso de que su documento reproduzca algún tipo de material que haya aparecido en una publicación anterior y cuyos derechos estén protegidos.

\subsection{Ecuaciones, fórmulas, matrices y términos matemáticos}

Deberán escribirse mediante algún editor de ecuaciones (por ejemplo, en \LaTeX, pero no como imágenes). Toda la simbología utilizada en ellas deberá estar claramente definida y distinguida. Las ecuaciones deberán justificarse a la izquierda y numerarse en forma consecutiva con números arábigos colocados a la derecha de la ecuación entre paréntesis. Se recomienda que las matrices se designen con letras mayúsculas resaltadas y los vectores con letras minúsculas resaltadas. Por ejemplo:

\begin{equation}
    x_{1,2} = \frac{-b \pm \sqrt{b^2 - 4ac}}{2a}
\end{equation}

\subsection{Tablas}

Para construir las tablas, utilice la opción incluida en su procesador de texto. 

\begin{table}[H]
    \centering
    \caption{Velocidades $VV$ de ondas compresionales a diferentes profundidades.}
    \begin{tabular}{r|ccccccc}
    \toprule
        $y_i$ [m] & 0 & 2 & 4 & 6 & 8 & 10 & 12 \\
        $vv(y_i)$ [m/s] & 290 & 690 & 840 & 900 & 930 & 940 & 950 \\
    \bottomrule
    \end{tabular}
    \label{T:tabla}
\end{table}

No repita información presente en el texto y utilice el formato TABLA I, TABLA II… para darles nombre. Según IEEE, debajo del nombre de “TABLA [número]”, se debe incluir un pequeño título que describa lo contenido en la tabla, el cual debe estar escrito en mayúscula. Se debe especificar la fuente de la tabla; si fue tomada de otra fuente se debe mencionar dicha fuente.

\begin{table}[H]
\caption{Otra tabla de ejemplo}
\label{T:otratabla}
\centering
\begin{tabular}{c l r r}
\toprule
\multicolumn{2}{c}{Producto} \\
\cmidrule(r){1-2}
Cantidad & Descripción & Precio unitario & Precio total  \\
\midrule
3  & Transistores 	& 250	& 750	\\
4  & Osciladores   	& 500   & 2000 \\
3  & Amp Ops     	& 600   & 1800	\\
10 & Resistores  	& 25    & 250	\\
10 & Capacitores	& 50 	& 500	\\
\midrule 
\multicolumn{3}{r}{TOTAL} & \textbf{5300} \\
\bottomrule
\end{tabular}
\end{table}

\subsection{Figuras}

Deberán seguir un estilo sencillo y claro. 

\begin{figure}[H]
    \centering
    \subcaptionbox{Primera figura \label{F:integrado}}
        {\includegraphics[width=0.2\textwidth]{integrado}}
    \subcaptionbox{Segunda figura \label{F:transistor}}
        {\includegraphics[width=0.2\textwidth]{transistor}}
    \caption{Ejemplo de dos figuras.}
    \label{animals}
\end{figure}

Al igual que en el caso de las tablas, no se debe repetir información presente en el texto. Por favor no incluya gráficos directamente extraídos de una hoja de datos y utilice un programa para darles una calidad más alta. Si fue tomada de otra fuente, esta se debe mencionar. Se debe indicar claramente las escalas y las unidades utilizadas, así como símbolos y abreviaturas empleadas. 

\begin{figure}[H]
    \centering
    \input{imagenes/histograma}
    \caption{Una gráfica de ejemplo}
    \label{F:ejemplo}
\end{figure}

Para citar las figuras, se utiliza la forma Fig. \ref{F:ejemplo} en el texto (evitar utilizar frases como ``la figura siguiente'' o similar pues su posición en el texto puede cambiar). En el caso de las subfiguras se citan como Fig. \ref{F:integrado} o, por ejemplo, la Fig. \ref{F:transistor}.

\subsection{Citas dentro del texto}

Las referencias aparecen en la misma oración de la referencia, entre paréntesis cuadrados y dentro de la puntuación correspondiente, por ejemplo: ``como muestra Marchette \cite{marchette2005random} y
como se mencionó anteriormente en
Horn \cite{horn1986robot}''. También las referencias entre paréntesis pueden funcionar como sustantivos: ``como se demostró en \cite{huang2009coupled} y de acuerdo con \cite{zamora2019smart}''.

Si desea citar las palabras exactas de algún autor se deben utilizar comillas (``'') y agregar el número de página dentro del paréntesis. Por ejemplo: 
Como resultado, esto ``crearía mejores espacios virtuales'' \cite[p. 56]{lafont2007protein}.

Si la cita tiene más de tres líneas, se recomienda separarla en un párrafo aparte con un tamaño de letra más pequeña. Por ejemplo:

Como menciona Wilde:

\begin{quote}
    \lipsum[1]
    \cite[p. 537]{zeng2014optimized}
\end{quote}

Si necesita citar más de una fuente en la misma oración, se citan las diferentes fuentes separadas por comas (,) o por un guion (-). Por ejemplo: 

\begin{itemize}
\itshape
    \item Varios trabajos anteriores han enfatizado el problema en cuestión \cite{marchette2005random}, \cite{horn1986robot}, \cite{huang2009coupled}
    \item Varios trabajos anteriores han enfatizado el problema en cuestión \cite{marchette2005random,horn1986robot,huang2009coupled}
\end{itemize}

Es importante destacar que IEEE prefiere evitar el uso de referencias citadas en otras referencias, es decir, de un autor citado en el trabajo de otro, por lo que aconseja que se localice la fuente original para ser citada.

\subsection{Sección de referencias}

Indicaciones generales:

\begin{itemize}
    \item Las referencias deben estar numeradas entre paréntesis cuadrados y alineadas a la izquierda. 
    \item En las referencias, los nombres de autores se abrevian a la inicial respectiva, precediendo el apellido. 
    \item Si la referencia tiene más de seis autores, se usa el nombre del primer autor y se le agrega el et al. Sin embargo, si la referencia contiene seis autores o menos, se debe mencionar a todos en la misma. 
    \item Todas las referencias deben terminar con un punto (.) al final, a menos que estas contengan un URL. 
    \item Las referencias no deben combinarse y sólo debe haber una referencia por número. 
    \item Deberá considerar que, si su artículo está escrito en inglés, este deberá seguir las normas de puntuación dadas en el manual de IEEE como, por ejemplo, el uso de la coma (,) dentro de las comillas para nombrar el título de un capítulo de libro o de un artículo. Por el contrario, si su artículo está escrito en español, deberá seguir las reglas de puntuación en este idioma, por ejemplo, el uso de la coma (,) fuera de las comillas para nombrar títulos. 
    \item Se recomienda el uso del DOI, en aquellas referencias que cuenten con este, por encima de la dirección URL de las fuentes en línea. 
    \item Para consultas más específicas, puede revisar \href{https://ieeeauthorcenter.ieee.org/wp-content/uploads/IEEE-Reference-Guide.pdf}{este} enlace de IEEE para el formato en inglés.
\end{itemize}

\begin{description}
    \item[Nota] Esta plantilla en \LaTeX\ tiene el programa de citas BibLaTeX en formato IEEE, el cual sigue automáticamente todo el formato descrito.
    \item[Bibliografías] En esta plantilla, todas las referencias bibliográficas se encuentran con el formato BibTeX en el archivo \texttt{referencias.bib}.
\end{description}

% Finalizar la numeración de líneas
\nolinenumbers

%%%%%%%%%%%%%%%%%%
\printbibliography
%%%%%%%%%%%%%%%%%%

%%%%%%%%%
\licencia

%%%%%%%%%%%%%%
\end{document}
%%%%%%%%%%%%%%